\section{Extensions to (un)typed term calculi}
\label{sec:term-calculi}
In light of the reduction to the negative fragment at the end of the previous section, we shall only consider types formed from $\N, \times, \arrow, \mu$ henceforth.


\subsection{From (co)derivations to (co)terms: rules as combinators}
It will be convenient for us to extend our computational model from just (co)derivations to a larger class of untyped (co)terms.
The main technical reason behind this is to allow the definition of a higher-order computability model necessary for our ultimate totality argument for $\cmuLJ$.
% In the presence of fixed points (or second-order quantifiers), such models require a domain of untyped terms or programs that are nonetheless computationally well behaved (giving rise to `candidates'), in order to handle the inherent impredicativity of second-order quantification in a non-circular way.
At the same time, we obtain a compressed notation for (co)derivations for notational convenience, and indeed carve out typed (conservative) extensions of the proof calculi thusfar considered.

In what follows, we use the metavariables $\rrule $ etc.\ to vary over inference steps of $\muLJneg$ and/or $\muLJnegnorec$, i.e.\ instances of any inference rules in these systems.
% For simplicity of exposition we shall henceforth construe $\muLJ$ (and $\muLJneg$) as the set of their inference steps, as only slight abuse of notation.
\begin{definition}
[(Co)terms~\cite{Das21:CT-fscd}]
 A \emph{coterm}, written $s,t$ etc., is generated coinductively by the grammar:
 \[
 s,t \ \bnf \ \rrule  \ \vert \ s\, t
 \]
I.e.\ coterms are possibly infinite expressions (of depth $\leq \omega$) generated by the grammar above.
A coterm is a \emph{term} if it is a finite expression, i.e.\ generated inductively from the grammar above.
 If all steps in a (co)term are from a system $\mathsf R$, we may refer to it as a $\mathsf R$-(co)term.
% 
% We write $\Coterms$ for the set of coterms and $\Terms$ for the set of terms. \anupam{do we need this? Maybe use minus superscript?} \todo{propose deletion}
 % We follow usual bracketing conventions, e.g.\ writing $stu$ for $((st)u)$.
\end{definition}

Our notion of (co)term is untyped, in that an application $st$ may be formed regardless of any underlying typing. 
(Co)terms will be equipped with a theory that (a) subsumes cut-reduction on (co)derivations; and (b) results in a computational model that is Turing complete.
Before that, however, let us see how (co)derivations can be seen as (co)terms.

\begin{definition}
    [(Co)derivations as (co)terms]
    We construe each $\cmuLJneg$ coderivation as a coterm (and each $\muLJneg$ derivation as a term) by identifying rule application with term application:
    % \begin{itemize}
    %     \item 
        if $\der$ ends with an inference step $\rrule$ with immediate sub-coderivations $\der_1 , \dots, \der_n$ 
        then $\der$ is $\rrule\, \der_1\, \dots \, \der_n$.
    % \end{itemize}
    
    Given a set $A$ of (co)terms, the \emph{closure} of $A$, written $\closure A$, is the smallest set of coterms containing $A$ and closed under application, i.e.\ if $s,t \in \closure A$ then also $st \in \closure A$.
\end{definition}

Of course if $\der$ is a derivation, then it is also a term.
Of particular interest to us in this work will be the class $\closure \cmuLJneg$, essentially finitary applications of progressing regular $\muLJnegnorec$-coderivations.


\begin{example}
    [Iterator coderivation as a regular coterm]
    \label{ex:iter-coder-as-coterm}
    Recalling the decomposition of the iterator as a circular coderivation in Example~\ref{ex:(co)iter-to-cycles}, let us specialise to the variation $    \vlinf{\iter{}}{}{\mu X \sigma(X)\seqar \tau}{\sigma(\tau)\seqar \tau}$.
%     We have that $\iter{}\, \der$ is expressed by the regular coderivation:
%     \[
%     \iter{}'(P):\quad 
%     \vlderivation{
% \vlin{\mulunf}{\bullet}{\mu X \sigma(X) \seqar \tau}{
% \vliin{\cut}{}{\sigma(\mu X \sigma(X))\seqar \tau}{
%     \vliq{\sigma}{}{\sigma(\mu X \sigma(X))\seqar \sigma(\tau)}{
%     \vlin{\mulunf}{\bullet}{\mu X \sigma(X)\seqar \tau}{\vlhy{\vdots}}
%     }
% }{
%     \vltr{P}{\sigma(\tau)\seqar \tau}{\vlhy{\ }}{\vlhy{}}{\vlhy{\ }}
%     % \vliin{\lr \arrow}{}{\sigma(\tau)\arrow\tau,\sigma(\tau)\seqar \tau}{
%     %     \vlin{\id}{}{\sigma(\tau)\seqar \sigma(\tau)}{\vlhy{}}
%     % }{
%     %     \vlin{\id}{}{\tau\seqar \tau}{\vlhy{}}
%     % }
% }
% }
% }
%     \]
 By following Example~\ref{ex:(co)iter-to-cycles}, we can express $\iter{}\, \der$
    by a regular coderivation, say $\iter{}'(\der)$, that, viewed as a coterm, satisfies the (syntactic) equation,
    \begin{equation}
        \label{eq:iter-coder-syntactic-eqn}
        \iter{}'(\der)\ = \ \mulunf\, (\cut\, \sigma(\iter{}'(\der))\, \der)
    \end{equation}
    % where $\sigma(\deri)$ is the functor coderivation from \todo{ref}. 
    Note that $\iter{}'(\der)$ above is indeed a \emph{regular} coterm: it has only finitely many distinct sub-coterms.
\end{example}

\anupam{commented `regular coterms' below}

% \begin{remark}[Regular coterms]\label{rem:reg-coterms}
% Just as we identified the regular coderivations as a natural finitarily representable subset of the coderivations, we may speak of the \emph{regular coterms} too, those with only finitely many distinct sub-coterms.
% For instance $\iter{}'(\der)$ in the example above is regular (as long as $\der$ is).
% Note that, since regular coterms are closed under application, we indeed have that $\closure \cmuLJ$ contains only regular coterms.

% Just as for regular coderivations \anupam{did we say this?} note that every regular coterm can be defined by a finite system of equations between coterms, e.g.\ as in \eqref{eq:iter-coder-syntactic-eqn} above.
% % of the form:
% % \[
% % \begin{array}{rcl}
% %   t_i & = &  \rrule_i\ t_1 \ldots  t_n\ \qquad  0 \leq i \leq n
% % \end{array}
% % \]
% % for some 
% % $t \dfn t_0, t_1, \ldots, t_n$ and for some rules $\rrule_0, \ldots, \rrule_n$. 
% % We will denote with $\Sys t$  the finite system of equations associated with $t$.  
% \end{remark}


% \begin{definition}[Farcs as systems of equations]\label{prop:system-of-regular-equations}
% For every $t\in \farcs$, we denote with $\Sys t$  the finite system of equations between farcs such that, for some 
% $t \dfn t_0, t_1, \ldots, t_n$ and for some rules $\rrule_0, \ldots, \rrule_n$, $\Sys t$ contains the following equations:

% since regular coterms have only finitely many distinct sub-coterms.
% \end{definition}


\subsection{Computational models: theories of (co)terms}
\anupam{need to make clear that we include `open' derivations among (co)terms too? Should check implications of this, probably just a matter of convenience.}

Let us henceforth make the following abbreviations:
\[
\pair \cdot \cdot : 
% \vlderivation{
% \vliin{\rr \times}{}{\sigma, \tau\seqar \sigma \times \tau}{
%     \vlin{\id}{}{\sigma \seqar \sigma}{\vlhy{}}
% }{
%     \vlin{\id}{}{\tau\seqar \tau}{\vlhy{}}
% }
% }
\vlderivation{
\vliin{\rr \times}{}{\seqar \sigma \times \tau}{
    \vlhy{\seqar \sigma}
}{
    \vlhy{\seqar \tau}
}
}
\qquad
\proj i:
\vlderivation{
\vlin{\lr \times}{}{\sigma_0 \times \sigma_1 \seqar \sigma_i}{
\vlin{\wk}{}{\sigma_0, \sigma_1 \seqar \sigma_i}{
\vlin{\id}{}{\sigma_i \seqar \sigma_i}{\vlhy{}}
}
}
}
\]
 % $\Niter$ for $\vliinf{\lr\N}{}{\N\seqar \sigma}{\seqar\sigma}{\sigma\seqar \sigma}$, $\iter{}$ for $\vlinf{\lr \mu}{}{\mu X \sigma(X)\seqar \tau}{\sigma(\tau)\seqar \tau}$,\anupam{need to expand these/explain}, $\injX\sigma $ for $\vlinf{\rr\mu}{}{\seqar \mu X\sigma(X)}{\seqar \sigma(\mu X \sigma(X))}$ and $\proj i$ for,
\[
% \begin{array}{r@{\ :\ }c}
    % \injX\sigma  & 
%     \vlderivation{
% \vlin{\rr\mu}{}{\sigma(\mu X \sigma(X))\seqar \mu X \sigma(X)}{
% \vlin{\id}{}{\sigma(\mu X \sigma(X))\seqar \sigma(\mu X \sigma(X))}{\vlhy{}}
% }
% }  
% \\
   % \iter{}  & 
   % \vlderivation{
   % \vlin{\lr \mu}{}{\sigma(\tau ) \arrow \tau , \mu X \sigma(X)\seqar \tau}{ 
   % \vliin{\lr \arrow}{}{\sigma(\tau) \arrow \tau, \sigma(\tau)\seqar \tau}{
   %  \vlin{\id}{}{\sigma(\tau)\seqar \sigma(\tau)}{\vlhy{}}
   % }{
   %  \vlin{\id}{}{\tau \seqar \tau}{\vlhy{}}
   % }
   % }
   % } 
   % \\
   % \Niter  & 
   % \vlderivation{
   % \vliiin{\lr \N}{}{\sigma, \sigma \arrow \sigma, \N \seqar \sigma}{ 
   %  \vlin{\wk}{}{\sigma, \sigma \arrow \sigma \seqar \sigma}{
   %  \vlin{\id}{}{\sigma \seqar \sigma}{\vlhy{}}
   %  }
   % }{
   %  \vlin{\wk}{}{\sigma, \sigma \arrow \sigma, \sigma \seqar \sigma}{
   %  \vliin{\lr \arrow}{}{\sigma \arrow \sigma, \sigma \seqar \sigma}{
   %      \vlin{\id}{}{\sigma \seqar \sigma}{\vlhy{}}
   %  }{
   %      \vlin{\id}{}{\sigma \seqar \sigma}{\vlhy{}}
   %  }
   %  }
   % }{
   %  \vlin{\id}{}{\sigma\seqar \sigma}{\vlhy{}}
   % }
   % }
% \end{array}
\injX\sigma : \vlinf{\rr\mu}{}{\seqar \mu X\sigma(X)}{\seqar \sigma(\mu X \sigma(X))}  
\qquad
\iter{\sigma} :  \vlinf{\lr \mu}{}{\mu X \sigma(X)\seqar \tau}{\sigma(\tau)\seqar \tau}
\]
\[
\numeral 0 : \Nzero \qquad \succ : \vlinf{\Nsucc}{}{\seqar \N}{\seqar \N}\qquad 
\Niter : 
   \vliinf{\lr\N}{}{\N\seqar \sigma}{\seqar\sigma}{\sigma\seqar \sigma}
\]
for $i \in \{0,1\}$.
We may omit types from subscripts when unimportant.
When writing (co)terms using the derivations above, we employ the convention that $\proj i$ and $\succ$ bind stronger than other applications, i.e.\ we write simply $t\, \proj i u$ for $t\, (\proj i u)$ and $t\, \succ u$ for $t\, (\succ\, u)$.

\todo{explain why we factor the $\times$ interactions (and even $\mu$ too) through other constants/gadgets: $\times$ to facilitate totality argument, and $\mu$ for similarity/to slightly facilitate realisability (though not really since we need bespoke iterators there).}
\todo{explain that we don't include $\inj$ for $\N$ to avoid $+$, instead this is hardcoded into the rules}
\anupam{recall that, apart from in, the actual types are not important due to untypedness.}

When referring to an arbitrary instance of a rule, the specification should be understood to be as originally typeset, unless otherwise indicated.
In particular, we follow this convention to define our notion of reduction on coterms:


\begin{figure}
    \centering
\[
\begin{array}{rcl}
  \ax \  x   &\rewrite  & x \\
  % \lr \unit \ t \ \vec x \ y  &\rewrite& t \ \vec x  \\
\ex \ t \ \vec x \ x \ y \ \vec y&\rewrite& t\  \vec x  \ y \ x \ \vec y  \\
\wk \ t \  \vec x\  x     & \rewrite & t\  \vec x\\
\contr \ t \ \vec x\ x &\rewrite&  t \ \vec x \ x \ x \\
\cut \ s \ t \ \vec x\ \vec y &\rewrite& t \ \vec y \ (s \ \vec x)
% \\
% \rr \imp \ t  \ \vec x \ x &\rewrite& t  \ \vec x \ x \\ 
\end{array}
\qquad
\begin{array}{rcl}
     {\rr \times \ s \ t \ \vec x \ \vec y }&\rewrite& \pair {  s\, \vec x}{t\,  \vec  y}
\\
 \lr \times \ t \, \vec x  \,y & \rewrite &  t\ \vec x \ \proj 0 y\ \proj 1 y 
 \\
 \proj i \pair {x_0} {x_1} & \rewrite & x_i \\
 \rr \imp \, t  \ \vec x \ x  &\rewrite&  t  \ \vec x \ x \\ 
\lr \imp \, s \ t  \ \vec x\ \vec y \ z  &\rewrite& t \ \vec y \ (z  \, (s \ \vec x )) 
%  \sem{\rr \mu} \ t \ \vec x &\betae&  t \ \vec x \\
%  \sem{\lr \mu} \ t \ \vec x &\betae&  t \ \vec x \\
%   \sem{\rr \nu} \ t \ \vec x &\betae&  t \ \vec x \\
%  \sem{\lr \nu} \ t \ \vec x &\betae&  t \ \vec x 
% \sem{\lr \mu} \ t \ \vec x \ (\sem{\rr \mu} \ u \ \vec x)  &=& t \ \vec x \ (A (\sem{ \lr \mu } )\ \vec x \ (u \ \vec y)) \\
% \sem{\lr \nu} \ t \ \vec x \ (\sem{\rr \nu} \ s \ \vec y \ z)&=& t \ \vec x\  (A(\sem{\rr \nu}) \ \vec y \ (s \ \vec y \ z)) 
\end{array}
\]
    \caption{Reduction for $\LJneg$ (both $\reduces$ and $\reducesnorec$). }
    \label{fig:reduction-ljneg}
\end{figure}


\begin{figure}
    \centering
    \[
\begin{array}{rcl}
 \rrule \ t \ \vec x &\reducesnorec&  t \ \vec x \qquad  \rrule\in \{\rr \mu, \mulunf \} 
\end{array}
\]
    \caption{Reduction for least fixed point rules in $\muLJnorec$. }
    \label{fig:reduction-mu-norec}
\end{figure}

\begin{figure}
    \[
    \begin{array}{r@{\ \reduces\ }l}
         \lr \mu \, s \, t\, \vec x\, \vec y\, z & t\, \vec y\, (\iter{}\, (s\, \vec x)\, z) \\
         \rr \mu\, t\, \vec x & \injX\sigma (t\, \vec x) \\
         \iter{}\, t\, x & x\, t \\
         \injX \sigma\,  x\, t & t\, (\sigma(\iter\, t)\, x)
         % \iter\, t\, (\injX\sigma\, x) & t\, (\sigma(\iter\, t)\, x)
    \end{array}
    \]
    \caption{Reduction for least fixed points in $\muLJ$.}
    \label{fig:reduction-mu-iter}
\end{figure}


\begin{figure}
  \[
   \begin{array}{r@{\ \reduces\ }l}
         \lr \N\, s\, t\, u\, \vec x\, \vec y\, z & u\, \vec y\, (\Niter\, (s\, \vec x)\, (t\, \vec x)\, z) \\
        { \Nsucc \, t \, \vec x  }&{ \succ \, (t\, \vec x) } \\
         \Niter\, s\, t\, \numeral 0 & s \\
         \Niter\, s\, t\, \succ x & t\, (\Niter \, s\, t\, x)
    \end{array}
  \]
    \caption{Reduction for $\N$ in $\muLJneg$}
    \label{fig:reduction-N-muLJ}
\end{figure}


\begin{figure}
   \[
   \begin{array}{r@{\ \reducesnorec\ }l}
        \Ncnd \, s\, t\, \vec x\, \numeral 0 & s\, \vec x \\
        \Ncnd \, s\, t\, \vec x\, \succ y & t\, \vec x\, y
   \end{array}
   \]
    \caption{Reduction for $\N$ in $\muLJnegnorec$.}
    \label{fig:reduction-N-muLJnorec}
\end{figure}


\begin{definition}
    [Theories]
    We define two (context-closed) reduction relations on (co)terms: 
    \begin{itemize}
        \item $\reduces$ is generated by the clauses in Figures~\ref{fig:reduction-ljneg}, \ref{fig:reduction-mu-iter} and \ref{fig:reduction-N-muLJ}.
        \item $\reducesnorec$ is generated by the rules in Figures~\ref{fig:reduction-ljneg}, \ref{fig:reduction-mu-norec} and \ref{fig:reduction-N-muLJnorec}. 
    \end{itemize}
    In all cases the lengths of vectors $\vec x$ and $\vec y$ match those of the relevant contexts $\Gamma$ and $\Delta$ from the original typesetting of the rules, i.e.\ from \Cref{fig:sequent-calculus-mulj,fig:co-iteration-rules,fig:unfolding-rules,fig:native-n-rules-mulj,fig:native-n-rules+cutreds-muljnorec}.
    Note that the use of both variables and term metavariables in these reduction rules is purely to aid parsing. All reduction rules are closed under substitution.
    
    When referring to (fragments of) $\closure \muLJneg$ as a computational model, we typically mean with respect to $\reduces$, and when referring to (fragments of) $\closure \cmuLJneg$ as a computational model, we typically mean with respect to $\reducesnorec$. {However we also consider a (weakly) extensional version of  $\convertsnorec$:}
    \begin{itemize}
        \item $\convertsnoreceta$ is the closure of $\convertsnorec$ under the rule   $\vlinf{\extensional}{}{ t\convertsnoreceta t'}{t\, x \convertsnoreceta  t'\, x}$.
    \end{itemize}
    Above $x$ must be a fresh variable, not a general (co)term.
\end{definition}

Admitting some extensionality is not necessary to reason about representability, since extensionality can be eliminated for low type levels, but simplifies some of the theorem statements.

\begin{example}
[Iteration equations]
The fundamental equation for iteration is indeed derivable by $\reduces$:
% \[
% \iter{} \, t\, (\injX \sigma \, x) \ =\ t\, (\sigma(\iter\, t)\, x)
% \]
% is indeed derivable using $\reduces$:
\[
\begin{array}{r@{\ \reduces\ }l}
     \iter{}\, P\, (\injX\sigma\, x) & \injX \sigma \, x\, P  \\
     & P\, (\sigma(\iter\, P)\, x)
\end{array}
\]
For $\reducesnorec$, recalling Example~\ref{ex:iter-coder-as-coterm} and using its notation,
% given some $P:\sigma(\tau) \seqar \tau$ let us write $\iter{}'(P)$ for the (progressing, regular) coderivation:
% \[
% \vlderivation{
% \vlin{\mulunf}{}{\mu X \sigma(X) \seqar \tau}{
% \vliin{\cut}{}{\sigma(\mu X \sigma(X))\seqar \tau}{
%     \vliq{\sigma}{}{\sigma(\mu X \sigma(X))\seqar \sigma(\tau)}{
%     \vlin{\mulunf}{\bullet}{\mu X \sigma(X)\seqar \tau}{\vlhy{\vdots}}
%     }
% }{
%     \vltr{P}{\sigma(\tau)\seqar \tau}{\vlhy{\ }}{\vlhy{}}{\vlhy{\ }}
%     % \vliin{\lr \arrow}{}{\sigma(\tau)\arrow\tau,\sigma(\tau)\seqar \tau}{
%     %     \vlin{\id}{}{\sigma(\tau)\seqar \sigma(\tau)}{\vlhy{}}
%     % }{
%     %     \vlin{\id}{}{\tau\seqar \tau}{\vlhy{}}
%     % }
% }
% }
% }
% \]
% Now 
we can simulate the iteration equation for $\iter\, P$ with $\iter{}'(P)$:
\[
\begin{array}{r@{\ \reducesnorec\ }ll}
     \iter{}'(P)\, (\injX\sigma\, x) & \iter{}'(P)\, x & \text{by $\injX\sigma$ reduction} \\
     & \cut\, \sigma(\iter{}'(P))\, P\, x & \text{by $\mulunf$ reduction} \\
     & P\, (\sigma(\iter{}'(P))\, x) & \text{by $\cut$ reduction}
\end{array}
\]
\anupam{point out that, at this juncture, $\reduces$ and $\reducesnorec$ do not agree/are not compatible, unlike $\cutred^* \subseteq \cutrednorec^*$? check this.}
\end{example}

\todo{need to mention clashing reductions and problem of confluence. we address this shortle by giving an interpretation in lambda terms.}

More importantly for us, our notion of extensional reduction on coterms subsumes that of cut-reduction on coderivations.
Since we have identified coderivations as coterms, we may state this rather succinctly, constituting the main result of this subsection:
\begin{theorem}
[Extensional reduction includes cut-reduction]
\label{thm:reduction-inludes-cut-elimination}
% Let $\der, \der'$ be coderivations of $\seqar \N$ in $\cmuLJneg$. If $\der \cutconvnorec\der'$ then $ \der \convertsnorec\der'$. 
${\cutrednorec}\subseteq {\convertsnoreceta}$.
% $\cutconvnorec \ \subseteq\  \extendedeq$.
\end{theorem}
% \begin{IEEEproof} [Proof idea]
% We show that, if $\der \cutrednorec \der'$ then, for some $n \geq 0$ sufficiently large, and for any $t_1, \ldots, t_n$, it holds that $\der\, t_1 \, \ldots \,  t_n\convertsnorec \der'\, t_1 \, \ldots \,  t_n $. We then conclude by repeatedly applying the rule $(\eta)$.
% % We prove a more general statement. If $\der$ is a coderivation of $\sigma_1, \ldots, \sigma _n \seqar \gamma$ and $\der \cutrednorec  \der'$ then, for any $t_1, \ldots, t_n \in \closure \cmuLJneg$, it holds that  $\der\, t_1 \, \ldots \,  t_n\convertsnorec \der'\, t_1 \, \ldots \,  t_n $
% % \[
% % \der\, t_1 \, \ldots \,  t_n\convertsnorec \der'\, t_1 \, \ldots \,  t_n 
% % \]
% Details are repored in~\Cref{app:proofs-of-sec-term-calcili}.
%  \gianluca{Anupam. Have a look at it.}  
% \end{IEEEproof}
\begin{proof}
    We show that, if $\der \cutrednorec \der'$ then, for some $n \geq 0$ sufficiently large, and for any $x_1, \ldots, x_n$, $\der\, x_1 \, \ldots \,  x_n\convertsnorec \der'\, x_1 \, \ldots \,  x_n $. We then conclude by repeatedly applying the rule $(\extensional)$. \anupam{would be easier to not have two distinct steps and just proceed directly.}
% \anupam{I changed the $\vec t$ to variables, it is important for the extensionality rule to make sense. shold change below too.}

% Let $\der$ be a coderivation of $\sigma_1, \ldots, \sigma _n \seqar \gamma$. If t $\der \cutrednorec  \der'$ then, for any $t_1, \ldots, t_n \in \closure \cmuLJneg$: 
% \[
% \der\, t_1 \, \ldots \,  t_n\convertsnorec \der'\, t_1 \, \ldots \,  t_n 
% \]

Suppose $\der \cutrednorec  \der'$. It suffices to consider the case where  the last rule of $\der$ is the cut rule rewritten by the  cut-reduction step. Indeed, the latter implies the general statement by appealing to the context closure of $ \reducesnorec$. We only consider some relevant cases. 

If $\der$ has the form,
\[
\vlderivation{
\vliin{\cut}{}{\Gamma, \Delta, \Sigma \seqar \gamma}
   {\vliin{\rr \times}{}{\Gamma,\Delta \seqar \sigma \times \tau}{\vldrs{\der_1}{\Gamma \seqar \sigma }}{\vldrs{\der_2}{\Delta \seqar \tau}}
   }
   {
\vlin{\lr \times}{}{\Sigma , \sigma \times \tau \seqar \gamma}{\vldrs{\der_3}{\Sigma, \sigma, \tau \seqar \gamma}}
   }
   }
\]
we have:
\[
\arraycolsep=2pt
\begin{array}{rcl}
\cut \, (\rr \times \, \der_1\, \der_2)\, (\lr \times\, \der_3) \, \vec x\, \vec y\, \vec z&\reducesnorec&
\lr \times \, \der_3 \, \vec z\, (\rr \times \, \der_1 \, \der _2\, \vec x\, \vec y)\\
& \reducesnorec & \lr \times \, \der_3 \, \vec z\, \pair{\der_1 \, \vec x}{\der_2\, \vec y}\\
% & \reducesnorec &  \der_3 \, \vec u\, (\proj 0{\pair{\der_1 \, \vec t}{\der_2\, \vec s}})\, (\proj 1 {\pair{\der_1 \, \vec t}{\der_2\, \vec s}})\\
& \reducesnorec^*  &   \der_3 \, \vec z\, (\der_1\, \vec x)\, (\der_2\, \vec y)\\ \\ 
\cut \, \der_2 \, (\cut \, \der_1\, \der_3)\, \vec x\, \vec y\, \vec z  &\reducesnorec & 
   (\cut \, \der_1 \, \der_3)\, \vec x\, \vec z\, (\der_2\, \vec y)\\
  & \reducesnorec  &    \der_3\, \vec z\, (\der_1 \, \vec x)\, (\der_2\, \vec y)
\end{array}
\]

If $\der$ has the form,
\[
\vlderivation{
\vliin{\cut}{}{\Gamma, \Delta, \Sigma \seqar \gamma}
   {
   \vlin{\rr \to}{}{\Gamma \seqar \sigma \to \tau}{\vldrs{\der_1}{\Gamma, \sigma \seqar \tau}}
   }{
   \vliin{\lr \to}{}{\Delta, \sigma \to \tau, \Sigma\seqar \gamma}{\vldrs{\der_2}{\Delta \seqar \sigma}}{\vldrs{\der_3}{\Sigma, \tau \seqar \gamma}}
   }
}
\]
we have:
% \[\arraycolsep=2pt
% \begin{array}{rcl}
% \cut \, (\rr \to \der_1) \, (\lr \to  \der_2\, \der_3) \, \vec x\, \vec y \, \vec z\, x^*&\reducesnorec& \lr \to  \der _2 \, \der_3 \, \vec y \, \vec z\, (\rr \to  \der_1\, \vec x)\, x^*\\
% &\reducesnorec &  \lr \to \, \der _2 \, \der_3 \, \vec y \, \vec z\, (\der_1\, \vec x)\, x^*\\   
% &\reducesnorec &
% \der_3 \, \vec z\, (\der_1 \, \vec x\, (\der_2 \, \vec y))\, x^* \\ \\
% \cut \, (\cut \, \der_2\, \der_1)\, \der_3\, \vec x\, \vec y \, \vec z\, x^*&\reducesnorec & \der_3\, \vec z\, (\cut \, \der_2\, \der_1)\, \vec x\, \vec y\\
% &\reducesnorec & 
% \der_3 \, \vec z\, (\der_1 \, \vec x\, (\der_2 \, \vec y))\, x^*
% \end{array}
% \]
\[\arraycolsep=2pt
\begin{array}{rcl}
\cut\, (\rr \to \der_1)\, (\lr \to  \der_2\, \der_3)\, \vec x\, \vec y\, \vec z & \reducesnorec & \lr \to  \der_2\, \der_3\, \vec y\, (\rr \to \der_1\, \vec x)\, \vec z \\
    &\reducesnorec & \der_3 \, \vec z \, (\rr \to \der_1 \, \vec x \, (\der_2\, \vec y)) \\
    &\reducesnorec & \der_3 \, \vec z \, ( \der_1 \, \vec x \, (\der_2\, \vec y)) 
    \\
    \\
\cut \, (\cut \, \der_2\, \der_1)\, \der_3\, \vec x\, \vec y \, \vec z &\reducesnorec & \der_3\, \vec z\, (\cut \, \der_2\, \der_1\, \vec x\, \vec y)\\
&\reducesnorec & 
\der_3 \, \vec z\, (\der_1 \, \vec x\, (\der_2 \, \vec y))
\end{array}
\]
% Notice that the  case above is the one requiring a sufficiently large number of applications $x_1 \, \ldots \,  x_n$, as we  crucially require an extra input $x^*$ to fire the reduction step $(\rr \to \der_1\, \vec x)\, x^* \reducesnorec (\der_1\, \vec x)\, x^*$. 


If $\der$ has the form,
\[
\vlderivation{
\vliin{\cut}{}{\Gamma, \Delta \seqar \gamma}
    {
    \vlin{\rr \mu}{}{\Gamma \seqar \mu X \sigma(X)}{\vldrs{\der_1}{\Gamma \seqar \sigma (\mu X \sigma(X))}}
    }{
    \vlin{\lr \mu ' }{}{\Delta, \mu X \sigma (X) \seqar \gamma}{\vldrs{\der_2}{\Delta, \sigma (\mu X \sigma(X))\seqar \gamma}}
    }
}
\]
we have:
\[
\arraycolsep=2pt
\begin{array}[b]{rcl}
\cut \, (\rr \mu \, \der_1)\, \lr \mu' \, \der_2\, \vec x\, \vec y&\reducesnorec &  (\lr \mu' \, \der_2)\, \vec y \, (\rr \mu \, \der_1\, \vec x)\\
&\reducesnorec^* &   \der_2 \, \vec y \, ( \der_1\, \vec x)\\ 
\noalign{\medskip}
\cut \, \der_1 \, \der_2 \, \vec x\, \vec y&\reducesnorec &   \der_2 \, \vec y \, ( \der_1\, \vec x) 
\end{array}
\qedhere
\]
\end{proof}





 \subsection{An embedding into $\lambda$-terms}
 While the significant technical development of this work involves `totality' arguments, e.g.\ in Section~\ref{sec:totality} showing that the representable partial functions of $\cmuLJneg$ (under $\convertsnoreceta$) are total, we better address \emph{determinism} too.
 
 As it stands, $\convertsnorec$ and $\convertsnoreceta$ may fire distinct reductions on the same (co)terms, so there is a priori no guarantee that the output of representable functions is unique.
However
% {In particular we show `uniqueness of the output' for the number-theoretic functions represented in any system by means of any rewriting relation we have introduced so far. The `totality' of such functions is already established thanks of the totality argument using the HR predicates.  
this can  be shown by  defining a straightfoward  interpretation of coterms of $\closure \cmuLJneg$ into the (untyped) $\lambda$-calculus. We shall prove that   the (untyped) $\lambda$-calculus, under $\beta\eta$-reduction, simulates $\closure\cmuLJneg$ under $\reducesnoreceta$. This simulation relies on the fact that we can express regular coterms as a finite system of equations, which are known to always have solutions in the (untyped) $\lambda$-calculus.
As the techniques are rather standard, we shall be quite succinct in the exposition.

% \begin{theorem}
%     The (untyped) $\lambda$-calculus, under $\beta\eta$-reduction, simulates $\closure\cmuLJneg$ under $\reducesnoreceta$.
% \end{theorem}
% This simulation relies on the fact that we can express regular coterms as a finite system of equations, which are known to always have solutions in the (untyped) $\lambda$-calculus.
% % (details are in Appendix~\ref{subsec:proofs-of-uniqueness}).

% We shall work with the untyped lambda calculus $\Lambda$ as our underlying programming language.

\anupam{I commented a remark below that seems completely out of place.}
% \begin{remark}
% Let us point out that it would not make any difference to take just the natural numbers as our underlying model (inducing a Kleene-style type structure).
% However the current exposition avoids the need for bureaucratic codings and uniformity issues. 
% \end{remark}

% Moreover,  can show that the relations on natural-numbers represented by $\closure \muLJneg$ are, indeed, functions. 







% We shall assume a set $\var$ of (`first-order' or `term') \emph{variables}, written $x,y$ etc.
% We shall use the same symbols for individual variables when working in arithmetic later.
\emph{$\lambda$-terms}, written $s,t$ etc., are generated as usual by:
\[
s,t \quad ::= \quad 
x \quad | \quad (t\, s) \quad | \quad \lambda x \,  t 
\]
We write $\Lambda$ for the set of all terms.
% For this (and later) notion(s) of term we follow usual bracketing conventions, e.g.\ writing $stu$ instead of $((st)u)$, and we assume that $\lambda $ binds as strongly as possible (unless a dot is used, e.g.\ $\lambda x.\, t\, s$ should be read $\lambda x \, (t\, s)$).
The notion of `free variable' is defined as expected, and we write $\fv t $ for the set of free variables of the term $t$.






\newcommand{\doublepair}[2]{\langle \! \langle {#1}, {#2}\rangle \! \rangle}
\newcommand{\doubleproj}[2]{\pi_{#1}\,  {#2}}
\newcommand{\godel}[1]{\lceil {#1} \rceil}

\begin{figure}[t]
    \centering
 \[
 \small
\begin{array}{rcl}
%    \identity& \dfn & \lambda x.x \\
% \letid M  N & \dfn & MN   \\
 \true&\dfn& \lambda x,y. x\\
      \false &\dfn& \lambda x,y. y\\
 \ifthen s t u   & \dfn & s \, t \, u \\
      \doublepair s t & \dfn & \lambda z.z\, s\, t\\
     % \project 0 M & \dfn & \enc{\true}  \\
     % \project 1 M & \dfn & \enc{\false}  \\
      % \letpair M x y N &\dfn & M(\lambda x. \lambda y. N)\\
      \doubleproj{0}{} &\dfn&  \lambda z. z \, \true \\
         \doubleproj{1}{} &\dfn&  \lambda z. z \, \false \\
  \godel {\numeral 0} & \dfn & \doublepair \false  {\lambda x.x}   \\
    \godel {\numeral{n+1}} &\dfn  & \doublepair \true {\godel n} \\
    \succ &\dfn & \lambda x. \doublepair \true   x \\
     \pred  &\dfn & \lambda x. \doubleproj{1}{x}  \\
    \cnd &\dfn &\lambda x,y,z. \ifthen{(\doubleproj 0 z)}{x}{y\, ({\pred \, z})}  
      % \inject 0 M &\dfn& \doublepair {\enc \true} M \\
      % \inject 1 M &\dfn& \doublepair {\enc \false} M 
      \end{array}
      \begin{array}{rcl}
\ifthen {\enc{\true}} s t & \betaeta & s\\
\ifthen {\enc{\false}} s t & \betaeta & t\\
 % \letpair   {\doublepair M N} xy P  & \betaeta&P[M/x, N/y]\\
 \doubleproj i {\doublepair {s_0} {s_1}} & \betaeta & s_i\\
 \pred \, {\godel {\numeral 0}} &\betaeta & \godel 0\\
 \pred\,  {(\succ \, {\godel {\numeral n}})} &\betaeta & \godel {\numeral n}\\
  \cnd \, {\godel {\numeral 0}} \, s \, t &\betaeta & s\\
 \cnd \, {(\succ \, {\godel {\numeral n}})}\, s \, t &\betaeta & t \, \godel{n}\\
 \end{array}
\]

    \caption{Macros for $\lambda$-terms  and corresponding reductions.}
    \label{fig:abbreviation}
\end{figure}







\begin{figure}[t]
    \centering
   \[
   \small
   \begin{array}{rcl}
       \sem{\ax}& \dfn& \lambda x.x \\
        \sem{\ex}&\dfn &\lambda t. \lambda \vec u. \lambda x. \lambda y. \lambda \vec v. t\, \vec u \,y\,x\, \vec v \\
  \sem{\cut}&\dfn &\lambda t. \lambda s. \lambda \vec{u}.\lambda \vec v. s\, \vec v \, (t \,\vec u)  \\ 
  \sem{\contr}&\dfn &\lambda t. \lambda \vec u. \lambda x. t\,\vec u \,x\,x\\ 
  \sem{\wk}&\dfn &\lambda t. \lambda \vec u. \lambda x. t \,\vec u  \\
   \sem{\rr \arrow}&\dfn& \lambda t. \lambda \vec u. \lambda x. t \,\vec u\, x\\  
   \sem{\lr \arrow}&\dfn& \lambda t. \lambda s. \lambda \vec u. \lambda \vec v. \lambda x.  s\,\vec v\,(x\,(t\,\vec u))\\
 %  \sem{\rr \unit}&\dfn & \identity \\
 % \sem{\lr \unit}&\dfn &\lambda t. \lambda \vec u. \lambda x. \letid x {t \vec u}   \\
   \end{array}
   \qquad
   \begin{array}{rcl}
         \sem{\rr \times }& \dfn&  \lambda t. \lambda s. \lambda \vec{u}.\lambda \vec v . \doublepair {t\, \vec u}{s \, \vec v} \\
  \sem{\lr \times}&\dfn &\lambda t. \lambda \vec u. \lambda x. {t\,  \vec u \, (\doubleproj0 x) \, (\doubleproj1 x)}  
 %  \\
 % \sem{\cev{\rr +}}&\dfn& \lambda t.\lambda \vec u. \inject 0{t \vec u}  \\
 %  \sem{\vec{\rr +}}&\dfn& \lambda t.\lambda \vec u.\inject 1{t \vec u}    \\
 %  \sem{\lr +}&\dfn& \lambda t. \lambda s. \lambda \vec u. \lambda x. \ifthen {(\enc \true x)}{t(\enc \false x )\vec u}{s(\enc \false x) \vec u}
\\
\sem{\rr \mu}&\dfn &\lambda t. \lambda \vec{u}. t\, \vec u \\
  \sem{\lr \mu}&\dfn& \lambda t. \lambda \vec{u}. \lambda z. t \, \vec u \, z  \\ 
  \sem{\Nzero}&\dfn & \godel{0}\\
  \sem{\Nsucc}&\dfn& \lambda t.\lambda \vec x. \succ \, {(t\, \vec x)}\\
  \sem{\Ncnd} &\dfn & \lambda t. \lambda s. \lambda \vec u. \lambda x. \cnd\, {x}\, {(t \, \vec u)}\, {(s\, \vec u)}
   \end{array}
   \]
    \caption{Translation of  $\cmuLJneg$ into $\lambda$-terms.}
    \label{fig:translation-of-rules}
\end{figure}




We work with a standard equational theory on $\lambda$-terms. 
$\betaeta$ is the smallest congruence on $\lambdaterms$ satisfying:
\[
\lambda x t\,  s \betae  t[s/x] \qquad \qquad \lambda x ( tx) =_\eta t \quad (x \not \in  \FV(t))
% \quad
% \lambda x (tx) \eqter t
\]
\Cref{fig:abbreviation} displays some macros for $\lambda$-terms (and the corresponding reduction rules) we will adopt in this subsection.
 We define an interpretation of coterms in $\closure \cmuLJneg$ into $\Lambda$.  
 % In what follows, for the sake of convenience we shall assume that $\closure \cmuLJneg$ is endowed with a denumerable set of (term) variables $x, y, z, \ldots$.  
 % 
  % \gianluca{Does $ \closure \cmuLJneg$  have term variables? I remember that FARCS has variables. If not we have to say that we can add them w.l.o.g.}
 % 
 We start with an interpretation of the basic inference rules: 

\todonew{A: I think we need to interpret regular coderivations first, and then extend to terms of $\closure \cmuLJneg$}

 \begin{defn}[Interpreting rules] To each inference step $\rrule$ of $\cmuLJneg$ we associate a $\lambda$-term  $\sem \rrule$ as shown in~\Cref{fig:translation-of-rules}. 
 % We extend $\sem{\_}$ to terms (i.e., finite coterms) by inductively setting $\sem{t\, s}\dfn \sem{t}\, \sem{s}$.
 \end{defn}
 
It is easy to see that this interpretation preserves equations from~\Cref{fig:reduction-ljneg},~\ref{fig:reduction-mu-norec}, and~\ref{fig:reduction-N-muLJnorec}:

 

 \begin{lemma} \label{lem:rewriting-into-beta}The following equations hold in $\Lambda$:
    \[
\begin{array}{rcl}
 \sem \ax \  x   &\betaeta  & x \\
  % \lr \unit \ t \ \vec x \ y  &\rewrite& t \ \vec x  \\
 \sem\ex \ t \ \vec x \ x \ y \ \vec y&\betaeta& t\  \vec x  \ y \ x \ \vec y  \\
 \sem\wk \ t \  \vec x\  x     & \betaeta & t\  \vec x\\
 \sem\contr \ t \ \vec x\ x &\betaeta&  t \ \vec x \ x \ x \\
 \sem\cut \ s \ t \ \vec x \ \vec y &\betaeta& t \ \vec y \ (s \ \vec x)\\
% \rr \imp \ t  \ \vec x \ x &\betaeta& t  \ \vec x \ x \\ 
{ \sem{ \rr \times} \ s \ t \ \vec x \ \vec y }&\betaeta& \doublepair {  s\, \vec x}{t\,  \vec  y}
\\
  \sem{\lr \times} \ t \, \vec x  \,y & \betaeta &  t\ \vec x \ (\sem{\proj 0} y)\ (\sem{\proj 1} y )
%  \\
% \sem{ \proj 0} \doublepair {x} {y} & \betaeta & x 
\end{array}
\qquad
\begin{array}{rcl}
 \sem{\proj i }\doublepair {x_0} {x_1} & \betaeta & x_i \\
      \sem{ \rr \imp} \ t  \ \vec x \ x  &\betaeta&  t  \ \vec x \ x \\ 
 \sem{\lr \imp} \ s \ t  \ \vec x\ \vec y \ z  &\betaeta& t \ \vec y \ (z  \, (s \ \vec x )) \\
%  \sem{\rr \mu} \ t \ \vec x &\betaeta&  t \ \vec x \\
%  \sem{\lr \mu} \ t \ \vec x &\betaeta&  t \ \vec x \\
%   \sem{\rr \nu} \ t \ \vec x &\betaeta&  t \ \vec x \\
%  \sem{\lr \nu} \ t \ \vec x &\betaeta&  t \ \vec x 
% \sem{\lr \mu} \ t \ \vec x \ (\sem{\rr \mu} \ u \ \vec x)  &=& t \ \vec x \ (A (\sem{ \lr \mu } )\ \vec x \ (u \ \vec y)) \\
% \sem{\lr \nu} \ t \ \vec x \ (\sem{\rr \nu} \ s \ \vec y \ z)&=& t \ \vec x\  (A(\sem{\rr \nu}) \ \vec y \ (s \ \vec y \ z)) 
 \sem{ \rr \mu} \ t \ \vec x &\betaeta&  t \ \vec x \\
  \sem{ \lr \mu} \ t \ \vec x &\betaeta&  t \ \vec x \\
  \sem{ \Ncnd}\, s\, t\, \vec x\, \godel{0} &\betaeta& s\, \vec x \\
      \sem{   \Ncnd} \, s\, t\, \vec x\, \succ y &\betaeta&  t\, \vec x\, y
\end{array}
\]
Also, if $\sem{t}\, x \betaeta \sem{t'}\, x$, for $x\notin \fv t$, then $\sem{t}\betaeta \sem{t'}$ by $\eta$-reduction in $\Lambda$.
\end{lemma}

 We now show how to extend  $\sem{\_}$ to regular coderivations, by noting that any such coterm can be described by a finite system of simultaneous equations. 
 From here we rely on a well-known result that such finite systems of equations always admit solutions in the untyped lambda calculus, with respect to $\betaeta$ (see, e.g., \cite{Hindley}).

 \begin{definition}
 [Interpreting regular coderivations]
     Consider a $\cmuLJneg$ coderivation $P$ with subcoderivations $\vec P = P_1, \dots , P_n$ where $P = P_1$.
 Suppose each $P_i$ is concluded by an inference step $\rrule_i$ with immediate subcoderivations $\vec P_i$ (a list of regular coderivations among $\vec P$).
 Write $\Sys \der$ for the system of equations $\{x_i = \sem{\rrule_i} \vec x_i\}_{i=1}^n$, where we set $\vec x_i \dfn x_{i_1}, \dots, x_{i_k}$ when $\vec P_i = P_{i_1},\dots, P_{i_k}$.
% 
We define $\sem P$ to be some/any solution to $x_1$ of $\Sys P$ in $\Lambda$, with respect to $\betaeta$.

 From here we extend the definition of $\sem\cdot $ to all coterms in $\closure \cmuLJneg$ inductively as expected, setting $\sem{ts}\dfn \sem t \sem s$.
 \end{definition}

\todonew{A: I just inlined the actual equations into a definition above and skipped the intermediate lemma here. I've glossed over the narrative too.}

%  \begin{lemma}\label{prop:system-of-equations} Consider a $\cmuLJneg$ coderivation with subcoderivations $\vec P = P_1, \dots , P_n$.
%  Suppose each $P_i$ is concluded by an inference step $\rrule_i$ with immediate subcoderivations $\vec P_i$ (a list of regular coderivations among $\vec P$)
%  There exist a system $\Sys P$ of simultaneous equations:
%  \begin{equation}
%      \label{eq:eqn-system}
%      \begin{array}{rcl}
%  x_1 &= &   \rrule_1 \, \vec x_1\\
% & \vdots &\\
%  x_n &= &   \rrule_n \, \vec x_n
% \end{array}
%  \end{equation}
%  such that:
%  \begin{itemize}
%      \item Each $ \rrule_i$ is an inference step of $\cmuLJneg$ and each $\vec x_i $ contains only variables among $x_1,\dots, x_n$.
%      % \item If $t \reducesnorec^* s$ then $x_1 \rightsquigarrow_{\mathrm{r'}\cup \Sys t} s$.
%       \item If $t \reducesnorec^* s$ then $x_1 \rightsquigarrow^*_{\mathrm{r'}\cup \Sys P} s$, where $\Sys P$ is construed as a rewrite system by orienting each equation of \eqref{eq:eqn-system} left-to-right. 
%  \end{itemize}
 
%  \end{lemma}
% \begin{proof}
% [Proof sketch]
% By definition of regular coderivation, $P$ has only finitely many distinct sub-coderivations, say $P_1,\dots, P_n$. We associate each variable $x_i$ to the coderivation $P_i$.
% If $P_i$ concludes with an inference step $\rrule$ with immediate sub-coderivations ${P_{i_1}}, \dots, {P_{i_k}}$ (in order), we duly set $\rrule_i\dfn \rrule$ and $\vec x_i \dfn x_{i_1}, \dots, x_{i_k}$.

% \todonew{A: shouldn't second item be stated rather as a simulation, extending $\vec x/\vec t$?}
% % setting $\rrule_i$ to be the inference step that concludes $P_i$ and $\vec x_1 $ to be $ x_{i_1} , \dots, x_{i_k} $ when $P_{i_1}$ 
% %     Follows from the fact that $t \in \closure \cmuLJneg$ is essentially obtained by finite applications of regular coterms, so $t$ has finitely many distinct coterms. 
% \end{proof}
 


% Now, let $t \in \closure \cmuLJneg$. By the above Lemma there is a system of simultaneous equations $\Sys t=\{ x_i = t_i(x_1, \ldots, x_n) \}$. Since $ t_i(x_1, \ldots, x_n)$ are finite terms, we have $\lambda$-terms $\sem{t_i}$. Then, we consider the following system of equations in $\Lambda$: 
% \[
% \begin{array}{rcl}
%  x_1 &\betaeta&   \sem{t_1(x_1, \ldots, x_n)}\\
% & \vdots\  &\\
%  x_n &\betaeta &   \sem{t_n(x_1, \ldots, x_n)}
% \end{array}
%  \]
% Now we may set $\sem{t}\dfn s_1$, where $s_1, \ldots, s_n\in \Lambda$ are solutions (respectively) of the  above system of simultaneous equations\footnote{These solutions always exist (see, e.g.,~\cite{Hindley}) but are not necessarily unique.}







Now, immediately from the definition of $\sem\cdot$ and \Cref{lem:rewriting-into-beta} we arrive at our intended interpretation:

\begin{proposition}\label{prop:cut-elimination-reduction} If $t \in \closure \cmuLJneg$ and $t \reducesnoreceta s$ then $\sem{t}\betaeta \sem{s}$. \anupam{added eta to the reduction (and removed $*$ for typesetting reasons). this is needed in the next corollary and follows from the results referenced in the proof.}
% {\ }
% \begin{itemize}
%     \item If $t \in \regularterms$ and $t \cutelim s$ then $\sem{t}\betar \sem{s}$.
%     \item If $t \in \terms$ and $t \cutelim s$ then $\sem{t}\betar \sem{s}$.
% \end{itemize}
\end{proposition}
% \begin{proof}
%     Straightforward by Lemma~\ref{lem:rewriting-into-beta} and Proposition~\ref{prop:system-of-equations}.
% \end{proof}



From here, by confluence of $\beta\eta$-reduction on $\lambda$-terms, we immediately have:
\begin{corollary}[Uniqueness]\label{thm:uniqueness} Let $t \in \closure \cmuLJneg$. If $\cod m \convertsnoreceta t \convertsnoreceta \cod{n}$ then $n=m$.
% {\ }
% \begin{itemize}
%     \item If $\der \cutelim^* \cod{n}$ and  $\der \cutelim^* \cod{m}$ in  $\muLJ$ and $\cmuLJ$, then $n=m$.
%     \item If $t \rewrite^* \cod{n}$ and  $t \rewrite^* \cod{m}$ in  $\terms$ and $\regularterms$, then $n=m$.
% \end{itemize}
\end{corollary}
% \begin{IEEEproof}[Proof idea]
%     We interpret the equational theory of $\closure \cmuLJneg$ (based on $\convertsnorec$) into the standard equational theory of (untyped) $\lambda$-terms and $\beta$-equivalence. We show that  the interpretation  preserves the equations of $\convertsnorec$, and we conclude by appealing to confluence for $\lambda$-terms. Details are reported in~\Cref{subsec:proofs-of-uniqueness}.
% \end{IEEEproof}
\begin{proof}
    Clearly, by Proposition~\ref{prop:cut-elimination-reduction} we have $\sem{t} \betaeta \sem{\cod{n}}$ and $\sem{t} \betaeta \sem{\cod{m}}$. Since $\sem{\cod{n}}\betaeta \godel{\cod n}$, we have $\sem{t} \betaeta \godel{\cod n}$ and $\sem{t} \betaeta \godel{\cod m}$. Since $ \godel{\cod n}$ and $\godel{\cod m}$ are normal forms, by confluence of  $\Lambda$ it must be that  $ \godel{\cod n}=\godel{\cod m}$, and hence $n=m$.
\end{proof}









\subsection{From typed terms back to proofs}\label{subsec:from-typed-terms-back-to-proofs.}
Let us restrict our attention to $\muLJneg$-terms in this subsection. 
{In what follows, for a list of types $\vec \sigma=(\sigma_1, \ldots, \sigma_n)$,  we write  $\vec \sigma \to \tau$ for $\sigma_1 \to \ldots \to \sigma_n \to \tau$}.
In order to more easily carry out our realisability argument in Section~\ref{sec:realisability}, it will be convenient to work with a typed version of $\clomuLJneg$:

% \todo{type assignment}
\begin{definition}
    [Type assignment]
    \emph{Type assignment} is the smallest (infix) relation `$:$' from terms to types satisfying:
    \begin{itemize}
        \item for each step $\vliiinf{\rrule}{}{\vec \sigma \seqar \tau}{\vec \sigma_1\seqar \tau_1}{\cdots}{\vec \sigma_n\seqar \tau_n}$ we have $\rrule : (\vec \sigma_1 \arrow \tau_1) \arrow \cdots \arrow (\vec \sigma_n \arrow \tau_n)\arrow \vec \sigma \arrow \tau$.
        \item if $t:\sigma\arrow \tau$ and $s:\sigma$ then $ts:\tau$.
        \item if $t:\mu X \sigma(X)$ and $s:\sigma(\tau)\arrow \tau$ then $ts:\tau$.
        % \item if $t: (\sigma \arrow X)\arrow X$, and $\sigma $ positive in $X$, then $t:\mu X \sigma$.
        % \anupam{do we need this last one? check realisability again, there it was apparently necessary for subject reduction}
    \end{itemize}
  % \red{(Type assignment) derivations will range over $\metader$.}
  We write $\typedclomuLJneg$ for the class of typed $\clomuLJneg$-terms.
\end{definition}

% Naturally, by inspection of the reduction rules of $\converts$, we have `subject reduction':
% \begin{proposition}
%     [Type preservation]\label{prop:type-preservation}
%     If $t: \tau$ and $t\reduces t'$ then $t':\tau$.
% \end{proposition}
% \begin{proof}
%     It suffices to consider the case where the reduction fires the outermost redex, by appealing to context closure for $\reduces$. This boils down to checking that the reductions in~\Cref{fig:reduction-ljneg},~\Cref{fig:reduction-mu-iter} and~\Cref{fig:reduction-N-muLJ} preserve types. The first group of reduction rules is routine, while the third one is derivable, so we will focus on the reduction rules in~\Cref{fig:reduction-mu-iter}. 

% Suppose $\lr \mu \, s \, t\, \vec x\, \vec y\, z: \tau$.
% This means that $s: \Gamma \to \sigma(\rho)\to \rho$, $t: \Delta \to \rho \to \tau$, $\vec x: \Gamma$, $\vec y: \Delta$ and $z: \mu X \sigma$. 
% Since $\iter{}: (\sigma (\rho)\to \rho)\to \mu X \sigma \to \rho $, $\iter{}\, (s\, \vec x): \mu X \sigma \to \rho$. Putting everything together, $t\, \vec y\, (\iter{}\, (s\, \vec x)\, z): \tau$ as required.

% For the redex of a $\rr\mu$ reduction, we must have  $\rr \mu\, t\, \vec x : \mu X \sigma $ with $t: \Gamma \to \sigma (\mu X \sigma )$ and $\vec x:\Gamma$. 
% So $t \, \vec x:\sigma (\mu X \sigma ) $. Hence, $\injX\sigma (t\, \vec x): \mu X \sigma  $.

% Suppose $\iter{}\, t\, x: \tau$. Then, $t: \sigma (\tau)\to \tau$ and $x: \mu X \sigma$. So $x\, t: \tau$.

% Finally suppose $\injX \sigma\,  x\, t: \tau$. Then  $x: \sigma (\mu X \sigma)$ and $t: \sigma (\tau)\to \tau$. This means that $\iter\, t: \mu X \sigma \to \tau$, and so $\sigma(\iter\, t): \sigma (\mu X \sigma) \to \sigma(\tau)$. Therefore,  $t\, (\sigma(\iter\, t)\, x): \tau$.


% \todonew{A: Here, how do we know that $\sigma(\iter\, t) : \sigma(\mu X \sigma) \to \sigma(\tau)$? It feels like we need a functoriality lemma here. Also, is this Proposition even important? It feels like (especially now that we don't include the construction clause for $\mu$s), the Theorem below can be established directly, with the local simulations of $\reduces$ by $\cutred$ inlining whatever type preservation is needed.}
% \end{proof}
% % \begin{IEEEproof}
% %     \red{See~\Cref{app:from-typed-terms-back-to-proofs.}.}
% % \end{IEEEproof}

The main result of this subsection is:
\begin{theorem}
    [Terms to derivations]\label{thm:terms-to-derivations}
  The natural number functions represented by $\typedclomuLJneg$ are already representable by $\muLJneg$.
\end{theorem}
\begin{proof}
[Proof sketch]
First, given a derivation $\der$ of ${\seqar \sigma \to \tau}$ in $\muLJneg$, we define the derivation  $\mathsf{uncurry}(\der)$ of $\sigma  \seqar \tau$ as  follows:
\[
\vlderivation{
\vliin{\cut}{}{\sigma \seqar  \tau}
{
\vldr{\der}{\seqar \sigma \to \tau}
}
{
\vliin{\lr \to}{}{\sigma, \sigma \to \tau \seqar \tau}{\vlin{\id}{}{\sigma \seqar \sigma}{\vlhy{}}}{\vlin{\id}{}{\tau \seqar \tau}{\vlhy{}}}
}
}
\]
We now define an interpretation of type assignments $t:\tau$ into derivations $\der_{\metader} $ of $ \seqar \tau$ by induction on $t:\tau$ as follows:
\begin{itemize}
     \item For each step $\rrule : (\vec \sigma_1 \arrow \tau_1) \arrow \cdots \arrow (\vec \sigma_n \arrow \tau_n)\arrow \vec \sigma \arrow \tau$, $\der_\rrule$ is the derivation of $\rrule$ in $\muLJneg$:
     \[
     \vlderivation{
     \vliq{\rr \to}{}{\seqar (\vec \sigma_1 \arrow \tau_1) \arrow \cdots \arrow (\vec \sigma_n \arrow \tau_n)\arrow \vec \sigma \arrow \tau}{
     \vliiin{\rrule}{}{\vec \sigma_1 \to \tau_1, \dots, \vec \sigma_n \to \tau_n , \vec \sigma \seqar \tau}{
        \vliq{}{}{\vec \sigma_1 \to \tau_1, \vec \sigma_1 \seqar \tau_1}{\vlhy{}}
     }{
        \vlhy{\cdots}
     }{
        \vliq{}{}{\vec \sigma_n \to \tau_n, \vec \sigma_n \seqar \tau_n}{\vlhy{}}
     }
     }
     }
     \]
    
    \item If $t:\sigma \to \tau$ and $s: \sigma$ then $\der_{ts}$ is defined as: 
    \[
    \vlderivation{
    \vliin{\cut}{}{\seqar \tau}{
        \vltr{\der_{s}}{\seqar \sigma}{\vlhy{\quad }}{\vlhy{}}{\vlhy{\quad }}
    }{
        \vltr{\colorbox{white}{$\scriptstyle\mathsf{uncurry}(\der_{t})$}}{\sigma \seqar \tau}{\vlhy{\quad }}{\vlhy{}}{\vlhy{\quad }}
    }
    }
    \]
    \item If $t: \mu X \sigma(X)$ and $s:\sigma(\tau) \to \tau$ then $\der_{ts}$ is defined as:
    \[
\vlderivation{
\vliin{\cut}{}{\seqar \tau}
{
\vldr{\der_{t}}{\seqar \mu X \sigma (X)}
}
{
\vliin{\lr \mu}{}{ \mu X \sigma (X)\seqar \tau}
{
    \vldr{\colorbox{white}{$\scriptstyle\mathsf{uncurry}(s)$}}{\sigma(\tau)\seqar  \tau}
    }
    {
    \vlin{\id}{}{\tau \seqar \tau}{\vlhy{}}
    }
    }
    }
    \]
\end{itemize}


\todonew{A: now that type assignment is deterministic, no longer need dependency on derivation. I've changed the proof reflecting that, old cases commented below.}
%  We now define an interpretation of typing derivations $\metader \proves t:\tau$ into (possibly open) derivations $\der_{\metader} $ of $ \seqar \tau$ by induction on $\metader$ as follows:
%         % \anupam{NB: $\der_\metader$ depends not only on $t$, but the typing derivation, and the decomposition $\vec \sigma \to \tau$.}
% \begin{itemize}
%     \item For each step $\rrule : (\vec \sigma_1 \arrow \tau_1) \arrow \cdots \arrow (\vec \sigma_n \arrow \tau_n)\arrow \vec \sigma \arrow \tau$, $\der_\rrule$ is the derivation of $\rrule$ in $\muLJneg$:
%     \[
% \vlderivation{
% \vldr{\der_\rrule}{\seqar (\vec \sigma_1 \arrow \tau_1) \arrow \cdots \arrow (\vec \sigma_n \arrow \tau_n)\arrow \vec \sigma \arrow \tau}
% }
%     \]

    
%     \item if $\metader\proves ts:\tau$ is inferred from $\metader_0 \proves t:\sigma \to \tau$ and $\metader_1 \proves s:\sigma$ then $\der_\metader$ is
%     \[
%     \vlderivation{
%     \vliin{\cut}{}{\seqar \tau}{
%         \vltr{\der_{\metader_1}}{\seqar \sigma}{\vlhy{\quad }}{\vlhy{}}{\vlhy{\quad }}
%     }{
%         \vltr{\mathsf{uncurry}(\der_{\metader_0})}{\sigma \seqar \tau}{\vlhy{\quad }}{\vlhy{}}{\vlhy{\quad }}
%     }
%     }
%     \]
%     \item if $\metader \proves ts:\tau$ is inferred from $\metader_0 \proves t:\mu X \sigma(X)$ and $\metader_1 \proves s: \sigma(\tau) \to \tau$ then 
%     % by induction hypothesis, we have:
%     % \[
%     % \vlderivation{
%     % \vldr{\der_{\metader_1}}{\seqar \sigma(\tau)\to \tau}
%     % }
%     % \]
% we define $\der_\metader$  as follows:
%     \[
% \vlderivation{
% \vliin{\cut}{}{\seqar \tau}
% {
% \vldr{\der_{\metader_0}}{\seqar \mu X. \sigma (X)}
% }
% {
% \vliin{\lr \mu}{}{ \mu X. \sigma (X)\seqar \tau}
% {
%     \vldr{\mathsf{uncurry}(\der_{\metader_1})}{\sigma(\tau)\seqar  \tau}
%     }
%     {
%     \vlin{\id}{}{\tau \seqar \tau}{\vlhy{}}
%     }
%     }
%     }
%     \]
%     % \anupam{I think we need to mention proof application here to somewhat justify the necessary equations.}
%     % \item if $\metader\proves t:\mu X \sigma(X)$ is inferred from $\metader'\proves t:(\sigma(X)\arrow X)\arrow X$ then, by induction hypothesis, we have a derivation $\der_{\metader'} (X) $ of $ \seqar (\sigma(X)\arrow X) \to X$, and so we have a derivation $\mathsf{uncurry}(\der_{\metader'} (X)) $ of $(\sigma(X)\arrow X)   \seqar  X$. We define $\der_{\metader}$ by applying a cut rule to the followiing derivations:
%     % \[
%     % \small
%     % \vlderivation{
%     % % \vliin{\cut}{}{\seqar \mu X \sigma(X)}{
%     %     \vlin{\rr \to}{}{ \seqar \sigma(\mu X \sigma(X)) \to   \mu X \sigma(X)}{\vlin{\rr \mu}{}{ \sigma(\mu X \sigma(X)) \seqar  \mu X \sigma(X)}{\vlin{\id }{}{\sigma(\mu X \sigma(X)) \seqar  \sigma(\mu X \sigma(X))}{\vlhy{}}}}
%     % % }{
%     % }
%     % \]
%     % \[ \small
%     % \vlderivation{
%     %     \vltr{\der_{\metader'}(\mu X \sigma(X))}{\sigma(\mu X \sigma(X))\arrow \mu X \sigma(X) \seqar \mu X \sigma(X)}{\vlhy{\quad}}{\vlhy{\ }}{\vlhy{\quad }}
%     % % }
%     % }
%     % \]
%     % where the substitution of $\mu X \sigma(X)$ for $X$ in $\der_{\metader'}(X)$ is justified by Remark~\ref{rem:gen-identity+substitution}.
% \end{itemize}
% % In all cases, we may need to use some further cuts against simple proofs depending on the precise decomposition $\vec \sigma \to \tau$.


Now, to show that  typed-$\closure \muLJneg$  and $\muLJneg$ represent the same functions on natural numbers, it suffices to prove that if $ t: \tau$  and $t \reduces t'$ then $\der_{t}\cutred^* \der_{t'}$.
% 
For this let us observe:
\begin{itemize}
    \item Whenever $t$ is typed, so are all its subterms by definition of type assignment.
    \item Any derivation $\der_{t(s)}$ in the form $\der_t(\der_s)$, for appropriate $\der_t(\cdot)$ (with leaf $\cdot$).
\end{itemize}
Thus the simulation of any reduction step $t(s) \reduces t(s')$, with redex $s$, is reduced to showing $s\reduces s' \implies \der_s \cutred^* \der_{s'}$.
% 
% Notice that  a type assignment derivation    $\metader' \vdash t:\tau$ exists   by Proposition~\ref{prop:type-preservation}. This requires to consider the case where the reduction fires the outermost redex, as  this allows us to conclude  by appealing to contextual closure for $\cutred$. 
This boils down to checking that the reductions in~\Cref{fig:reduction-ljneg},~\Cref{fig:reduction-mu-iter} and~\Cref{fig:reduction-N-muLJ} are simulated by a series of cut-reduction rules on $\muLJneg$, which  is routine. 
% \end{proof}
% 
    % [Proof idea]
    %  We define an interpretation of typing derivations $\metader \proves t : \tau$ into (possibly open) derivations $\der_{\metader}$ of   $\seqar \tau$ by induction on $\metader$.
%     Define an interpretation of typing derivations $\metader \proves t:\vec \sigma \to \tau$ by (possibly open) derivations $\der_{\metader} : \vec \sigma \seqar \tau$ by induction on $\metader$ as follows:\anupam{NB: $\der_\metader$ depends not only on $t$, but the typing derivation, and the decomposition $\vec \sigma \to \tau$.}
% \begin{itemize}
%     \item For each step $\rrule$, $\der_\rrule$ is just itself, $\rrule$.
%     \item if $\metader\proves ts:\tau$ is inferred from $\metader_0 \proves t:\sigma \to \tau$ and $\metader_1 \proves s:\sigma$ then $\der_\metader$ is
%     \[
%     \vlderivation{
%     \vliin{\cut}{}{\seqar \tau}{
%         \vltr{\der_{\metader_0}}{\seqar \sigma}{\vlhy{\quad }}{\vlhy{}}{\vlhy{\quad }}
%     }{
%         \vltr{\der_{\metader_1}}{\sigma \seqar \tau}{\vlhy{\quad }}{\vlhy{}}{\vlhy{\quad }}
%     }
%     }
%     \]
%     \item if $\metader \proves ts:\tau$ is inferred from $\metader_0 \proves t:\mu X \sigma(X)$ and $\metader_1 \proves s: \sigma(\tau) \to \tau$ then
%     $\der_\metader$ is just the application $\der_{\metader_0}\der_{\metader_1}$, cf.~
%     Definition~\ref{def:appl-and-cut-red-muLJ}.
%     \item if $\metader\proves t:\mu X \sigma(X)$ is inferred from $\metader'\proves t:(\sigma(X)\arrow X)\arrow X$, with $\der_{\metader'} (X) : \sigma(X)\arrow X \seqar X $ obtained by the inductive hypothesis,
%     $\der_\metader$ is,
%     \[
%     \small
%     \vlderivation{
%     \vliin{\cut}{}{\seqar \mu X \sigma(X)}{
%         \vliq{\injX\sigma}{}{\seqar \sigma(\mu X \sigma(X)) \arrow \mu X \sigma(X)}{\vlhy{}}
%     }{
%         \vltr{\der_{\metader'}(\mu X \sigma(X))}{\sigma(\mu X \sigma(X))\arrow \mu X \sigma(X) \seqar \mu X \sigma(X)}{\vlhy{\quad}}{\vlhy{\ }}{\vlhy{\quad }}
%     }
%     }
%     \]
%     where the substitution of $\mu X \sigma(X)$ for $X$ in $\der_{\metader'}(X)$ is justified by Remark~\ref{rem:gen-identity+substitution}.
% \end{itemize}
% In all cases, we may need to use some further cuts against simple proofs depending on the precise decomposition $\vec \sigma \to \tau$.
% It remains to check the simulation of cut-reductions.\todo{Gianluca} 
\end{proof}

% \todo{from typed $\closure \muLJ$ back into $\muLJ$}





